\section{Institutional Background and Related Literature}
\label{sec:background}

This section situates India's December 2021 suspension of agricultural derivatives
within the country's long history of futures bans, summarises the unsettled
theoretical and empirical debate over whether futures trading destabilises spot
markets, sets out the minimum-support-price (MSP) and procurement apparatus that
confounds any naive before--after comparison, and locates our contribution against
the most directly competing studies. Full citations are collected in
\texttt{references.md}; throughout we use author--year keys, and we flag conflicts
of interest where the evidence is industry-commissioned.

\subsection{A recurrent Indian reflex: futures suspensions, 2007--2021}
\label{sec:background-history}

Suspending or delisting agricultural futures in response to food-price inflation is
a recurring feature of Indian commodity policy rather than a novel intervention. The
institutional lineage runs through a sequence of government committees---Dantwala
(1966), Khusro (1980), Kabra (1994) and the Abhijit Sen Expert Committee
(2008)---each convened amid inflation anxiety and each cautious about (re)permitting
futures in politically sensitive staples. The Dantwala Committee warned as early as
1966 that suppressing futures prices irrespective of spot behaviour destroys the very
utility of a futures market, anticipating the identification problem at the heart of
the present evaluation.

The modern episode begins with the 2003 reintroduction of exchange-traded
commodity futures, followed almost immediately by a wave of delistings: in February
2007 wheat and rice were delisted, and in 2008 further commodities (including chana,
urad, tur, potato, soya oil and rubber) were suspended amid the global food-price
spike \citep{Nath2008, Babshetti2019}. Gulati, Chatterjee and Hussain (2017)
document roughly fifteen suspensions between 2003 and 2017 and show that suspension
duration correlates with a commodity's weight in the food basket---direct evidence
that the policy responds to inflation salience rather than to any market-specific
diagnostic. A distinct, manipulation-driven episode is the March 2012 guar-seed and
guar-gum delisting, where the Forward Markets Commission documented genuine
over-speculation and a price bubble \citep{Singh2012, Soni2013}; this is the one
suspension in the lineage with a clean manipulation trigger, in contrast to the
diffuse inflation rationale of 2007--08 and 2021. We treat the manipulation-triggered
(2012 guar) versus inflation-triggered (2007--08, 2021) distinction as analytically
important, because it bears on whether the suspended commodities were selected on a
transitory price shock.

The 2021 episode unfolded in stages. The Securities and Exchange Board of India
(SEBI), having absorbed the Forward Markets Commission in 2015, first directed the
National Commodity and Derivatives Exchange (NCDEX) to bar new positions in chana on
16 August 2021 and in mustard seed and oil on 8 October 2021. On 19--20 December 2021
SEBI suspended derivatives trading for one year across seven commodities---wheat,
paddy (non-basmati rice), the soybean complex, crude palm oil (CPO), and
moong---while continuing the chana and mustard suspensions; existing positions were
permitted only to square off. The suspension has since been extended six times and
now runs through 31 March 2027, making it a more-than-five-year halt with no published
ex-ante rationale and currently under SEBI review. The staggered onset (chana in
August, mustard in October, the remaining five in December) is consequential for
identification and motivates our use of commodity-specific treatment dates rather
than a single pooled event.

A further institutional fact shapes the cost side of the question: the banned
commodities accounted for the overwhelming majority of NCDEX agricultural turnover,
so the suspension effectively hollowed out the domestic agri-derivatives market and
left hedgers without listed instruments \citep{Rajib2024, PwC2023}. We note that the
two reports documenting this turnover collapse are NCDEX-funded and treat their
qualitative claims, not their causal magnitudes, as evidence.

\subsection{Do futures destabilise spot prices? Theory is inconclusive}
\label{sec:background-theory}

The premise underlying a derivatives suspension---that futures trading inflates or
destabilises the spot price---has no unconditional theoretical warrant. On the
stabilising side, Working (1949, 1960) and Gray (1963), exploiting the 1958 U.S.
onion futures ban, find cash-price variability generally \emph{higher} in years
without futures; Danthine (1978) shows in a rational-expectations setting that
futures provide forecasts guiding production and storage and so tend to increase spot
stability; and Cox (1976) argues that futures raise the information content of spot
prices, so that removing them degrades discovery. The ``onion experiment'' is the
canonical natural experiment for the proposition that removing futures \emph{raises}
spot volatility, and it frames our prior.

Against this, equilibrium storage models (Turnovsky 1979; Newbery 1987) establish
that futures can either stabilise or destabilise spot prices depending on parameters,
so that the sign of the effect is an empirical, not a theoretical, question. The
financialisation literature is similarly split: Irwin and Sanders (2011, 2012) find
little systematic evidence that index-trader positions move returns or volatility,
whereas Tang and Xiong (2012) document increased co-movement and volatility for
index-member commodities after 2004; Cheng and Xiong (2014) synthesise this as
operating through risk-sharing and information channels rather than a simple
speculation-inflates-prices mechanism. Critically for our design, Gupta and Varma
(2016) find for Indian rubber that spot volatility is \emph{both} a cause and a
consequence of trading activity---the precise simultaneity that makes any naive
comparison of banned commodities before and after suspension hazardous. We therefore
treat the +8--10\% ``volatility rose'' prior we inherited not as a theoretical
prediction to confirm but as a claim to stress-test against endogeneity-robust
evidence.

\subsection{Empirical evidence on Indian suspensions}
\label{sec:background-empirics}

The Indian evidence base is large but methodologically uneven, and---as our results
show---its headline conclusions are sensitive to whether a counterfactual is
constructed.

\paragraph{The naive before--after tradition.} A long line of within-India studies
compares volatility (or inflation) across pre-futures, with-futures and post-ban
regimes for the same commodities. Nath and Lingareddy (2008) report higher volatility
in pulses during futures trading; Sobti (2020), using an augmented E-GARCH with ban
and relaunch dummies on the 2007--08 episode, finds spot volatility highest during the
ban phase for most commodities; and, closest to our own outcome, Dey and Gairola
(2024) report that spot volatility \emph{increased} post-ban for chana and refined soy
oil after the 2021 suspension. These regime-comparison designs select treatment timing
on the outcome---commodities are banned \emph{because} their prices ran hot---and so
are vulnerable to the mean-reversion artifact (Ashenfelter's dip; Chab\'{e}-Ferret
2015) that our placebo battery is built to detect. We cite this tradition as the
literature we are stress-testing, not as a benchmark to match. Consistent with the
artifact interpretation, Sendhil and Ramasundaram (2014) find wheat spot volatility
\emph{low} during the 2007 ban precisely because it had spiked beforehand and then
mean-reverted.

\paragraph{The counterfactual tradition.} The most important methodological precedent
is the synthetic-control work of Aggarwal, Chatterjee and Sehgal (2022, 2023), who
build synthetic counterfactuals for chana (June 2016 and August 2021) and mustard
(October 2021) from a donor pool of non-suspended pulses and oilseeds, with predictors
spanning international prices, production and import shocks, and mandi establishment.
They find that the synthetic counterfactual tracks actual prices in every episode: the
suspension neither raised nor lowered \emph{price levels} or inflation. Their donor
pool, predictor set and per-episode (staggered) treatment map directly onto our
control-commodity problem, and their null is the result we extend to new margins. The
gap we fill is explicit: their estimand is the price \emph{level}, whereas our primary
outcome is spot \emph{realised volatility} (the H1/C1 margin), estimated on a
disaggregated district panel with modern few-treated inference.

The industry-commissioned reports reinforce this contrast. Gaurav and Pandey (2024),
in an NCDEX-commissioned study, report that a naive pre/post comparison shows
volatility up for all banned commodities, yet their own difference-in-differences of
volatility against domestic and international peer commodities is \emph{null} except
for soy oil at the twelve-month horizon---a within-study illustration that the
``volatility rose'' headline is an uncontrolled-comparison artifact. We adopt their
peer-mapping and the null DiD as corroboration while flagging the NCDEX funding.
Rajib, Barai and Arora (2024), also NCDEX-funded, measure ``volatility'' as
cross-mandi price dispersion---a spatial estimand distinct from our time-series
realised volatility---and document that this dispersion was already rising from April
2021, well before the December suspension, which supports the policy-endogeneity
reading of the ban's timing.

Looking back to the 2007--08 episodes, the weight of peer-reviewed Indian evidence is
that suspensions did not curb inflation of the banned commodities (Babshetti and
Basanna 2019; Sahoo and Kumar 2009; Singh and Goyal 2011), and the government's own
Sen Committee (2008) could not establish a causal link from futures trading to
agricultural price inflation. Bose (2008) further finds that Indian agricultural
futures indices were informationally weak even before the bans, which tempers how much
price-discovery loss any suspension can be expected to cause.

\subsection{MSP and procurement as a confounder}
\label{sec:background-msp}

A first-order threat to any volatility comparison in this setting is administered
pricing. For commodities subject to a minimum support price backed by Food
Corporation of India or state-agency procurement, the realised price is pinned at or
near the support level for long stretches, so measured ``volatility'' reflects the
procurement regime rather than market conditions. This is not a marginal concern in
our sample: paddy's daily spot returns are exactly flat 40.3\% of the time, against
1--6\% for clean (non-administered) commodities, because procurement holds the price at
the support level; wheat is similarly though less severely censored at 19.7\% flat
returns. Realised volatility computed on a government-administered price is a
mechanical artifact, not market volatility. We therefore drop paddy from the primary
volatility analysis as effectively censored, and retain wheat---one of our core
treated commodities---only with an MSP flag, reading it as a robustness check rather
than a clean estimate. More broadly, the MSP/procurement apparatus is a confounder
that must be carried explicitly in the design (it enters our confounder ledger) rather
than assumed away, and it interacts with the contemporaneous export bans, stock limits
and edible-oil duty changes that overlapped the suspension window for several of the
banned commodities.

\subsection{Where this paper sits}
\label{sec:background-gap}

This paper occupies a specific and, to our knowledge, unfilled position in the
literature. First, on \emph{estimand}: the existing counterfactual work measures price
levels and inflation, while the naive within-India literature measures volatility
without a counterfactual. We measure spot \emph{realised volatility} against an
explicitly constructed counterfactual, combining the strengths of both traditions. The
inherited claim that the suspension raised volatility by roughly 8--10\% does not
survive. In the current food-donor rerun, the headline two-way fixed-effects
difference-in-differences on the district panel puts spot realised volatility
$9.8\%$ lower, not higher, than its counterfactual, but the aggregate estimate is
not statistically significant (CR1 $p \approx 0.145$; wild-cluster bootstrap
$p \approx 0.153$). Synthetic difference-in-differences (Arkhangelsky et al. 2021)
is larger in magnitude, at roughly $-19.1\%$, but remains borderline
($z \approx -1.88$). The commodity-level signal is concentrated in chana:
synthetic DiD estimates about $-38.7\%$ and the Abadie synthetic-control estimate is
about $-37.3\%$, with placebo inference at the attainable floor of the current
small donor set. Wheat is near zero in the current SDID run and remains heavily
confounded by MSP procurement and contemporaneous export-policy interventions. We
therefore frame the result as a \emph{suggestive}, chana-centered quasi-causal
association rather than a precise aggregate effect: the suspension is associated
with lower spot volatility in the treated basket, but the pooled food-donor
estimate is modest and statistically weak.

Second, on \emph{method}: because the design has very few treated clusters (five to
six treated commodities), a single adoption window, and endogenous treatment timing,
we draw on the synthetic-control family \citep{Abadie2010, Arkhangelsky2021} and on
few-treated inference (wild-cluster bootstrap, conformal and Conley--Taber procedures)
rather than relying on conventional cluster-robust standard errors, which are known to
be unreliable below roughly thirty clusters \citep{MacKinnonWebb,
BertrandDufloMullainathan2004}. We also document a methodological correction that
proved decisive: an earlier version of the design, run on a calendar-grid spot series
in which mandi data behaved as seven-day grids while realised volatility was annualised
on a 252-trading-day convention, \emph{failed} its falsification tests (a fake-ban
placebo at $p=0.046$ and a pre-trend test at $p=0.033$). After a trading-day weekday
filter, the dropping of MSP-censored paddy, and the food-donor expansion, those
diagnostics are materially cleaner: the fake-December-2019 placebo is about
$-6.5\%$ and insignificant (analytic $p \approx 0.344$; wild bootstrap
$p \approx 0.433$), and the joint pre-trend test does not reject parallel trends
($p \approx 0.445$). This sequence is itself a caution about the fragility of the
naive designs that dominate the prior literature.

Third, on the companion volatility-mechanism (C2) margin, we find no support for the
``elevated volatility'' claim through a different lens: an ICSS-then-within-regime
procedure with the HAC-corrected Sans\'{o}--Arag\'{o}--Carri\'{o}n detector locates no
unconditional-variance break near the suspension date, consistent with the structural-
break literature's warning that an unmodelled level shift inflates apparent GARCH
persistence \citep{Hillebrand, Lamoureux1990}. The basis arm of the original design
is internally infeasible as stated---there are no post-ban futures for the banned
commodities, so no post-ban basis exists---and is accordingly reframed as a
domestic-versus-international spread question for future work, pending vendor futures
data. The honest characterisation of the contribution is therefore a both-arms
evaluation, not advocacy: we test whether the suspension delivered the price stability
it was implicitly meant to provide, and find that, if anything, banned-commodity spot
volatility was lower---not higher---under the suspension, while being explicit about
the limitations (few treated clusters, MSP censoring, the absence of post-ban futures,
and donor-pool construction still in progress).
